\documentclass{ximera}
\title{Pullbacks}

\begin{document}

    \begin{abstract} Pullbacks: interaction with  wedge product and exterior derivative.    \end{abstract}
    \maketitle

    A reference for this material is Chapters 12 and  14 of John M. Lee. Introduction to smooth manifolds. Second edition. Grad. Texts in Math., 218. Springer, New York, 2013. xvi+708pp. ISBN: 978-1-4419-9981-8.

    \begin{definition}[Pullbacks]
        Let $M$, $N$ be smooth manifolds, $f : M \to N$ a smooth function, and $\omega \in \Omega ^k (N)$. The \emph{pullback} of $\omega$ via $f$ is the $k $-form $f^{\ast} \omega \in \Omega ^{k}(M)$ given by 
        \[  f^{\ast} \omega (X_1, \ldots , X_k) : = \omega (  f_{\ast} X_1, \ldots , f_{\ast} X_k  ).          \]
    \end{definition}

    \begin{proposition}[We love pullbacks]
        The pullback behaves well with respect to wedge product and exterior derivative. Namely, if $\omega \in \Omega ^k (N)$ and $\eta \in \Omega ^{\ell}(N)$, then
        \[    f^{\ast} (\omega \wedge \eta ) = (f^{\ast} \omega ) \wedge (f^{\ast} \eta)                \]
        and 
        \[    d f^{\ast} \omega = f^{\ast } d \omega    .      \]
    \end{proposition}

    \begin{solution}
        For vector fields $X_1, \ldots , X_{k + \ell} \in \mathfrak{X}(M)$, we compute
\begin{align*}
     f^{\ast} (\omega & \wedge \eta )  (X_1, \ldots , X_{k + \ell } )  =  \omega \wedge \eta (f_{\ast} X_1 , \ldots , f_{\ast} X_{k + \ell } ) \\
     & = \frac{1}{k! \, \ell ! }\sum_{\sigma \in S_{k + \ell }} \omega (  f_{\ast} X_{\sigma (1) }, \ldots , f_{\ast} X_{\sigma (k)}  ) \eta (  f_{\ast} X_{\sigma (k + 1 )}, \ldots , f_{\ast} X_{\sigma (k + \ell )}  ) \\
     & = \frac{1}{k! \, \ell ! }\sum_{\sigma \in S_{k + \ell }}  f^{\ast}  \omega (   X_{\sigma (1) }, \ldots ,  X_{\sigma (k)}  ) f^{\ast} \eta (  X_{\sigma (k + 1 )}, \ldots , X_{\sigma (k + \ell )}  ) \\
     & = ( f^{\ast} \omega ) \wedge (f^{\ast} \eta ) (X_1, \ldots , X_{k + \ell}). 
\end{align*}
to prove the second identity, we consider a form 
\[     \omega = g \, dy^{i_1}\wedge \ldots \wedge dy ^{i_k} .     \]





        
    \end{solution}

    

\end{document}